% =====================================================================
%  Section 4 - Problem formulation (author's 4.1-4.2, edited; 4.3-4.4 in the same voice)
% =====================================================================
\section{Problem Formulation}
\label{sec:problem}

\subsection{Search Space and Two Tuning Stages}
\label{sec:space}

A candidate solution is one reel set $x = (r_1, \dots, r_5)$, where each strip $r_j \in \Sigma^{L}$ holds $L = 50$ symbols from the 9 symbols $\Sigma = \{\mathrm{A}, \dots, \mathrm{G}, \mathrm{W}, \mathrm{S}\}$. We put no constraint on symbol counts, so a strip can hold any number of copies of a symbol, or none.

The base and free reels are tuned separately, first the base reels and then the free reels. During the base stage, the three base stats (base RTP, hit rate and trigger rate) are calculated and tuned. For the free stage the base reels are fixed and only the free reel set changes. The three free stats are free RTP, free hit rate and retrigger rate. Free RTP also depends on the base trigger rate $T$, but keeping the base reels constant makes it depend only on the free reels.

\subsection{Objective}
\label{sec:objective}

Each stage has three stats $s_i(x)$, each with a target $t_i$ and a weight $w_i$ (Table~\ref{tab:targets}), and we minimise the weighted sum of relative errors
\begin{equation}
  F(x) \;=\; \sum_{i} w_i \, \frac{\lvert s_i(x) - t_i \rvert}{t_i}\,.
  \label{eq:fitness}
\end{equation}
$F(x) = 0$ means every target stat is matched exactly. Dividing the error of each stat by its target normalises it, so that stats on very different scales, such as an RTP of 0.57 and a trigger rate of 80, can be added together. Rates enter in their ``1 in $X$'' form. RTP has the highest weight, since it is the value that testing labs and licensing bodies check, whereas the trigger rate has more room: being 15 or more spins away from the target can be acceptable. A weighted sum is the simplest way to combine several targets and is close to how a mathematician reads a PAR sheet, but it hides the trade-off when targets conflict, which happens in the free stage (Sect.~\ref{sec:results}).

\begin{table}[t]
\centering
\caption{Targets $t_i$ and, in brackets, weights $w_i$ of the objective~\eqref{eq:fitness}. Rates are ``1 in $X$''; the free stage has a retrigger rate.}
\label{tab:targets}
\begin{tabular}{@{}lccc@{}}
\toprule
Stage & RTP & Hit rate & (Re)trigger rate \\
\midrule
Base & 0.57 (10) & 3 (6)   & 80 (3) \\
Free & 0.38 (10) & 2.7 (6) & 60 (3) \\
\bottomrule
\end{tabular}
\end{table}

\subsection{Noisy Evaluation by Simulation}
\label{sec:noise}

For this game, every stat could actually be calculated exactly: line-win RTP follows from the symbol counts on each reel, hit rate from going through all $50^5$ windows, and the free game from a geometric series over the number of free spins. However, we want a tuning process that also works for cascade and cluster games, where exact maths becomes much harder. Hence, the same way mathematicians check reel sets, we treat the game as a black box and use a multithreaded C++ Monte Carlo simulator. One evaluation plays $n = 10^7$ base spins with a random seed $\omega$ and returns $\hat F_n(x;\omega)$, which takes about 0.15\,s on 10 threads. In the free stage, the trigger rate $T$ of the fixed base reels is measured once from $10^8$ spins, and each evaluation then plays $\mathrm{round}(10^7/T)$ free rounds directly, i.e.\ the number of rounds a $10^7$-spin full game would play on average, and free RTP is calculated by dividing by $T$. This skips the base spins between free rounds and saves simulation time.

Since every stat comes from a finite number of spins, the score is noisy. Every evaluation in a run uses the same seed $\omega$ (common random numbers~\cite{law2015simulation}), so the same reel set always gets the same score, and a reel set that got a lucky score keeps it. The standard deviation of one $10^7$-spin estimate of $F$ is around 0.013 for good base reel sets and 0.035 in the free stage (Sect.~\ref{sec:noisestudy}), which near the end of the base stage is close to the gap between good candidates. A run also reports the lowest of about 1000 noisy scores, which favours the reel set whose noise was lucky~\cite{jin2005evolutionary,rakshit2017noisy}. So for every result we report, we re-evaluated the run's best reel set with $10^8$ spins on a seed the run never used (Sect.~\ref{sec:reeval}).

\subsection{Initialisation and Variation Operators}
\label{sec:operators}

A random reel set puts one copy of every symbol on each strip, fills the other 41 positions at random and shuffles the strip. Mutation picks $m \in \{1, 3, 5, 10\}$ of the 250 positions without replacement and changes each one to a different random symbol. The GA, hill climbing and simulated annealing all use this same mutation, so any difference between them comes from the population, crossover and the acceptance rule. Crossover works reel by reel: for each reel a cut point $c \in \{1, \dots, L-1\}$ is drawn, and the child takes the first $c$ symbols from one parent and the rest from the other. Every child gets crossover first and then mutation.
